\documentclass[journal]{IEEEtran}
\usepackage{cite}
\usepackage{amsfonts}
\usepackage{amssymb}
\usepackage{amsmath}
\usepackage{amsthm}
\usepackage{url}
\usepackage{enumitem}
\usepackage{array}
\usepackage{booktabs}
\usepackage[utf8]{inputenc}
\usepackage[english]{babel}

\newtheorem{definition}{Definition}
\newtheorem{assumption}{Assumption}
\newtheorem{lemma}{Lemma}
\newtheorem{law}{Law}
\newtheorem{theorem}{Theorem}
\newtheorem{proposition}{Proposition}
\newtheorem{obligation}{Obligation}

\newcommand{\Rg}{\mathcal{R}}
\newcommand{\E}{\mathcal{E}}
\newcommand{\Ok}{\mathsf{Ok}}
\newcommand{\Err}{\mathsf{Err}}
\newcommand{\Wait}{\mathsf{Wait}}
\newcommand{\Exhausted}{\mathsf{Exhausted}}
\newcommand{\Pool}{\mathcal{P}}
\newcommand{\use}{\operatorname{use}}
\newcommand{\liab}{\operatorname{liab}}
\newcommand{\issue}[1]{\#{#1}}

\author{Ryan~J.~Kung%
\thanks{Ryan J. Kung can be reached at ryankung@ieee.org.
Manuscript created October 25, 2026.  This document specifies the load bounds of a Rings node
before 1.0: what it holds node-wide, per neighbour edge and per lane, what a free identity can
consume, and what refusal and backpressure look like at each bound (\texttt{rings-transport},
module \texttt{callback}; \texttt{rings-core}, modules \texttt{fair\_admission},
\texttt{swarm::callback::inbound}, \texttt{chunk}, \texttt{message::replay} and
\texttt{dht::chord::carrier\_store}).  It is the law for issues \issue{907}, \issue{927},
\issue{932}, \issue{934}, \issue{938} and \issue{1094} of \texttt{RingsNetwork/rings}, and closes
the Sybil residual of the removal-stability law~\cite{RemovalStability}.}}
\title{\huge Load Bounds:\\
Node, Edge and Lane Accounting\\
for a Rings Node under Free Identities}
\date{}

\IEEEaftertitletext{%
\vspace{0.5\baselineskip}%
\noindent\hfill
\parbox{0.90\textwidth}{%
{\centering\bfseries Abstract\par}%
\vspace{0.4\baselineskip}%
\noindent
A Rings identity is a key digest, so anyone can mint as many as they like.  Every table a
remote party can grow must therefore be accounted to something that is not free.  We take the
authenticated neighbour edge: a connection the node itself admitted, which holds one of its
$\Rg$ connection slots.  One accounting law covers every such table: each slot has a reserved
floor, overflow shares a pool in fair claims, and the node-wide bound is the sum.  An honest
edge is then refused only past its slot's floor and its fair claim, and a free identity
consumes at most the share of the slot that carries it.  We apply the law to the
receiver's replay streams (with an eviction horizon that cannot reopen replay), to the removal
store, to received frames through a credit-reclaim protocol that makes the frame bound hard and
independent of the connection count, and to reassembled messages, which are charged once at
their first chunk so that no honest message is refused.  We state the complexity obligations of
the admission hot paths, and replace the reassembly barrier's control exemption by a barrier
per message category, under which order is guaranteed per category only.
}\hfill\null\par
\vspace{0.5\baselineskip}}

\begin{document}
\maketitle
\emergencystretch=1.5em

\section{Problem}\label{sec:problem}

\begin{description}[leftmargin=1.2em,style=nextline]
\item[\issue{907}] The receiver's replay table never evicts, and a new stream fails closed once
the table is full.  Streams are keyed by the origin account, which costs nothing, so 16\,384 key
pairs and one message each lock a node out of every new origin, permanently.
\item[\issue{1060}, residual] The removal store charges each sender a share, and senders are
free keys, so 64 identities fill it.
\item[\issue{934}] Credit bounds the frames held per connection, not per node: about 1.3\,GiB
at the connection cap.  Advertised credit cannot be taken back, so no hard node-wide bound
exists today.
\item[\issue{932}] A reassembled message is re-charged at its full size on the inbound actor,
where it cannot wait, and is dropped when that charge does not fit.  A chunk that finds the
reassembly budget full is dropped too.
\item[\issue{938}] Admission costs $O(n)$ per event and takes node-global locks per frame,
exactly when the node is saturated.
\item[\issue{927}] The reassembly barrier exempts control, which reorders chunked control, and
orders data across categories, which no protocol needs.
\end{description}

\section{Model}\label{sec:model}

\begin{definition}[Principals]\label{def:principals}
\begin{itemize}[leftmargin=1.2em]
\item An \emph{origin} $o$ is an account DID, recovered from a signature.  Anyone can create
one, at no cost.
\item An \emph{edge} $e$ is a handshake-authenticated connection the node admitted
(\texttt{EdgeRelation::Neighbour}).  The registry holds at most
$\Rg=2(F+s+1)$ connection records in any phase, where $F=160$ is the number of finger slots and
$s$ the number of successors ($328$ at $s=3$).  So at most $\Rg$ edges are \emph{current} at
any instant.  An edge identity is also a free key; what is bounded is the number of edges
current at once.
\item A \emph{slot} $\kappa\in K=\{0,\dots,\Rg-1\}$ is one unit of the registry's capacity.  An
edge takes the \emph{cleanest} free slot $\kappa(e)$ when it is admitted and frees it when it
is retired, both in the transition that announces it, so the current edges hold distinct slots.
A slot is not a topology reference: one edge may back several fingers and a successor, or none.
Let $\use^{t}_{T,\kappa}$ be the items of table $T$ (Def.~\ref{def:table}) charged to $\kappa$
that are still live at $t$: an item past its instant is not counted, since the next admission of
$T$ retires it.  The \emph{floor use} of $\kappa$ at $t$ is its worst table's,
$w_t(\kappa)=\max_T\min(\use^{t}_{T,\kappa},f_T)/f_T$, the part of a floor that an edge admitted
into $\kappa$ would inherit; the cleanest free slot is the one of least $w_t$, the least index
among equals.
\item A \emph{lane} $(e,\ell)$ is one of the four credit-controlled channels of an edge.
\item The \emph{local} principal $\lambda$ is the node itself: its own operations alone (its
own writes and repairs, and out-of-band connection signaling it is handed).
\item A \emph{dedicated} principal $d\in D_T$ is a duty of the node that remote input drives
but no slot introduces, with a share of its own in table $T$: the inbox principal $\mu$, which
introduces deliveries out of the node's own relay inbox (replay table only), and the custody
principal $\chi$, which introduces the hand-off of a range the node owns (removal store only).
\end{itemize}
\end{definition}

\begin{definition}[Table]\label{def:table}
A table $T$ is a node-wide set of items that remote input can grow: replay streams, removal
records, frames held, reassembly bytes, mailbox bytes.  Every item $x\in T$ carries a principal
$\pi(x)\in K\cup\{\lambda\}\cup D_T$: the slot of the edge that \emph{introduced} it, the local
principal, or a dedicated one, and it
leaves $T$ only by a transition of $T$'s own law (consumption, retirement at a horizon, or an
acknowledged move).
\end{definition}

\begin{assumption}[What a remote peer can send]\label{as:remote}
A remote party can:
\begin{itemize}[leftmargin=1.2em]
\item mint origins without bound;
\item hold up to $\Rg$ edges at once, and replace an edge by a fresh identity whenever
connection admission accepts one;
\item send, on each edge, anything the edge's credit, quotas and admission laws let through.
\end{itemize}
It cannot authenticate as an edge it does not hold.  Sybil and eclipse resistance remain
non-goals; this law bounds what a Sybil consumes, and does not prevent one.
\end{assumption}

\textbf{Refusal and backpressure.}  At a bound, a transition takes one of three forms:
\begin{itemize}[leftmargin=1.2em]
\item \emph{backpressure}: the transition waits for capacity in resource order, and nothing is
lost;
\item \emph{refusal}: it fails with an error that names the bound and changes nothing, and the
sender keeps its copy;
\item \emph{drop}: it is discarded and counted.
\end{itemize}
A drop is allowed only where the law names it (best-effort relay, the application inbox).  It
is never allowed for an honest frame or message at a capacity bound.

\section{The accounting law}\label{sec:accounting}

\begin{law}[Slot accounting]\label{law:edge}
Every table $T$ has a floor $f_T$ per slot, a pool $\Pool_T$, a local share $L_T$ and a share
$S_d$ per dedicated principal $d\in D_T$, with node-wide bound
\begin{gather*}
B_T(s)=\Rg(s)\, f_T+\Pool_T+L_T+\textstyle\sum_{d\in D_T}S_d,\\
\Rg(s)=2(F+s+1),
\end{gather*}
so a deployment that raises the successor count $s$ (\issue{774}, \issue{777}) raises the floors
it reserves, $2f_T$ per successor, and nothing else.  Let $\use_\kappa$ be the number of items
charged to slot $\kappa$, and $o_\kappa=\max(0,\use_\kappa-f_T)$ its pool items.  The
\emph{claimants} are the slots with $o_\kappa>0$ or a pending item over their floor; let
$n\ge1$ be their number.  An item introduced by $\pi$ is admitted iff:
\begin{itemize}[leftmargin=1.2em]
\item $\pi=\kappa$ and $\use_\kappa<f_T$ (the slot's floor); or
\item the pool has room, $\sum o<\Pool_T$, and the claimant is within its fair claim,
$o_\kappa<\lceil\Pool_T/n\rceil$; or
\item $\pi=\lambda$ and $\use_\lambda<L_T$; or
\item $\pi=d\in D_T$ and $\use_d<S_d$.
\end{itemize}
A transition that adds several items is admitted iff each of them is, in order, and is refused
whole otherwise.  A refused transition is refused or, where the law says so, it waits.  An item stays charged to its slot
until it leaves by $T$'s own law, whoever holds the slot: a departure changes no charge, and an
edge admitted into a slot inherits its usage.  No item is ever dropped to make room: a claimant
above its fair claim after $n$ grew keeps what it holds and gains nothing until it falls below.
\end{law}

\begin{theorem}[Isolation]\label{theorem:isolation}
Under Law~\ref{law:edge}:
\begin{enumerate}[leftmargin=1.5em]
\item $|T|\le B_T$ always;
\item an item through slot $\kappa$ with $\use_\kappa<f_T$ is never refused, however many
origins, edges or slots any other party controls; in particular an edge keeps its floor for as
long as it holds its slot, less what the slot held when it was admitted.  The guarantee is per
item: a transition adding $k$ items through $\kappa$ is admitted only if
$k\le(f_T-\use_\kappa)+(\lceil\Pool_T/n\rceil-o_\kappa)$ with the pool's room, and is refused
whole otherwise;
\item an over-floor slot is refused only once it holds its fair claim or the pool is full, and
it is guaranteed $f_T+\lceil\Pool_T/n\rceil\ge f_T+\lceil\Pool_T/\Rg\rceil$ items as the other
claimants' items expire, with $n-1$ of them;
\item the items of the slots a party holds or has held number at most $j f_T+\Pool_T$ for $j$
such slots, and at most $\Rg f_T+\Pool_T$ however many edges it rotates.
\end{enumerate}
\end{theorem}
\begin{proof}
(1)~No transition moves an item between slots or into the pool: a departure changes no charge.
Each slot holds at most $f_T$ items within its floor, the pool admits only while $\sum o<\Pool_T$,
the local principal holds at most $L_T$, and each dedicated principal at most $S_d$.
(2)~The floor test reads $\use_\kappa$ alone, and only transitions through $\kappa$ change it.
(3)~A claimant below $\lceil\Pool_T/n\rceil$ is refused only with a full pool, which then holds
some claimant above its claim; that claimant gains nothing and its items leave by $T$'s law, so
the pool drains to the fair claims, with $n\le\Rg$.
(4)~By (1), with $j\le\Rg$.
\end{proof}

\begin{proposition}[Rotation residual]\label{prop:rotation}
A party that rotates edges through free slots charges each slot it passes with up to $f_T$
items, which stay until they leave by $T$'s law.  An edge admitted at $t$ inherits
$\min_{\kappa\in F_t}w_t(\kappa)$ of its floors, over the free slots $F_t$, $|F_t|=\Rg-c$ with
$c$ edges current (Def.~\ref{def:principals}).  So:
\begin{enumerate}[leftmargin=1.5em]
\item \emph{Targeted.}  The party takes any of a newcomer's floor only if it has charged every
free slot with items live at the newcomer's admission: $\Rg-c$ admissions within the life $h$ of
those items, $h\le H+2\sigma$, or $H+H_I+3\sigma$ for application streams.
\item \emph{Concurrent.}  A party that holds $\Rg-c$ edges at once, fills their floors and
departs charges every free slot, under cleanest-first as under least-free, at the same $\Rg-c$
admissions per horizon.  Against it the choice of slot changes nothing.
\item \emph{Least-free.}  Had an edge taken the least free slot instead, (1) would fail: a
departing edge usually frees the least free slot, which the next admission takes, so one
admission per victim suffices.
\item An edge that holds its slot throughout is never affected.
\end{enumerate}
\end{proposition}
\begin{proof}
Law~\ref{law:edge}: a slot's charges outlive its edge.  (1)~If some free slot has $w_t=0$, the
newcomer takes one, and $w_t(\kappa)>0$ needs an item of $\kappa$ live at $t$, charged by an
edge admitted into $\kappa$ within its life.  (2)~With every free slot charged, every choice
inherits.  (3)~After the party's edge departs from the least free slot $\kappa$, $\kappa$ is the
least free slot unless a lesser one was freed meanwhile.  (4)~Theorem~\ref{theorem:isolation}(2).
\end{proof}

Cleanest-first is a strict improvement and not a new bound.  From any state, an admission
inherits no more under it than under least-free, since it takes the least $w_t$ over the same
free slots; and it removes the targeted attack of (3), one admission per victim, leaving only the
concurrent bound of (2).  Proposition~\ref{prop:rotation}(2) is the cost this law does not
remove: free keys and connection churn can leave newly admitted edges without a floor for one
horizon of $T$ per round of $\Rg-c$ admissions.  Such a party is
held only to its fair claim, $\lceil\Pool_T/n\rceil$, on the slots it charged, and an honest
newcomer keeps that claim: degradation, never a refusal of an existing edge, and never a
permanent lockout (Prop.~\ref{prop:nolockout}).

\textbf{Why not an admission rate limit (\issue{1094}).}  Keeping every free slot charged costs
$(\Rg-c)/h$ admissions per unit time, about $3$ a minute at $h=100$\,min, $\Rg=328$ and $c=28$,
and any positive rate at the application horizon.  An honest node admits new edges at $c/\bar L$,
with $\bar L$ the mean edge life: at most the mean session, $4.2$\,min for Weibull sessions of
shape $0.5$ and median $1$\,min with 90\% browser nodes~\cite{DurableRecords}, so about $7$ a
minute for a browser with $c\approx28$ and $70$ for a seed with $c\approx300$.  A limit that never
defers honest churn therefore admits more than the attack needs, and a limit tight enough to
bound it, burst plus refill below $\Rg-c$ per horizon, would defer a seed's restart (about $\Rg$
reconnections at once) by tens of minutes.  Keying it by address is rejected (\issue{780}; many
honest browsers share one address behind carrier-grade NAT), and keying it by slot is void, since
the node assigns slots.  An honest joining storm meets none of this: every arrival is admitted at
once into the cleanest slot.  The levers that remain are a cost on identity, which this law does
not assume, and position randomization (\issue{780}, 2.0), which takes from a fresh identity the
choice of where it lands.

Charging the edge identity instead of its slot is unsound: a departing identity's floor items
stay, while a fresh identity brings a fresh floor, so rotation alone grows $|T|$ past $B_T$.

Accounting to the origin, as the removal store did~\cite{RemovalStability} and the replay
table's quota does, fails part~(4).  An origin costs nothing, so a party controls unboundedly
many, and the only bound left is the node-wide total.  The slot is the unit that is both
authenticated and bounded in number.  The quota per origin stays: it is a rate within an edge,
not a capacity.

\textbf{Who overflows honestly.}  In an overlay of mostly browser nodes, the honest overflow is
on relay edges: a browser reached through a seed sees every origin the seed forwards to it on
that one edge.  Without the fair claim, a party with $k$ edges could take the whole pool and
refuse every honest relay edge for the table's horizon.  With it, a relay edge keeps
$\lceil\Pool_T/n\rceil$ however many slots the party holds.  The alternative of charging only
streams whose origin is the delivering edge, and leaving relayed streams to the origin quota,
reopens \issue{907}: a relayed final-destination stream is exactly the stream a free identity
creates, and the quota is per origin.

\section{Replay streams (\issue{907})}\label{sec:replay}

A receiver stream $s$ admits a transaction on two paths.  Live, only while its proof is live,
$t\le ts+ttl$, with $ttl\le\mathtt{MAX\_TTL\_MS}=H$ and $ts\le t+\sigma$.  Out of a relay inbox,
as of its hold instant $h$, which the witness bounds by the sender's proof, $h\le ts+ttl$; only
an application message is ever held.

\begin{definition}[Relay hold horizon]\label{def:relayhorizon}
A relay element held at $h$ retires at $h+H_I$, $H_I=\mathtt{MAX\_RELAY\_INBOX\_TTL\_MS}$
(7\,days): every replica projects it out of its carrier at a read, a join or a sync, and a drain
never delivers it.  The predicate reads the signed hold instant, not the element's dot, so it is
element-intrinsic and commutes with join, and a member that rewrites an old hold under a fresh
dot gains nothing.  An element is projected out, never refused, so a carrier with an old
element stays admissible.  Relay removes keep their cap.
\end{definition}

\begin{definition}[Horizon]\label{def:streamhorizon}
Let $x_s=\max(ts+ttl)$ over the transactions $s$ admitted, and
$\iota(\mathsf{Application})=H_I+\sigma$, $\iota=0$ for every other class.  The stream is
\emph{spent} at $t$ iff $t\ge x_s+\iota+\sigma$.  So a control, storage or E2E stream of
default-lifetime traffic is spent about $10$\,min after its last transaction, and none later
than $H+2\sigma$; an application stream $H_I+\sigma$ later.
\end{definition}

\begin{lemma}[Eviction is sound]\label{lem:evict}
A spent stream can be forgotten without reopening replay.  Every transaction it ever admitted is
then undeliverable at this node's clock, even one stepped back by $\sigma$, on either path, so a
stream opened afresh at $\mathsf{First}$ never sees one again.
\end{lemma}
\begin{proof}
Live, a transaction is refused at its proof once $t>ts+ttl$, and $t-\sigma\ge x_s$.  Out of an
inbox, it is delivered only while $t<h+H_I\le ts+ttl+H_I$, and $t-\sigma\ge x_s+H_I+\sigma$.
Both checks run before replay, whatever the window holds.
\end{proof}

\begin{law}[Stream admission]\label{law:streams}
The receiver table is a table under Law~\ref{law:edge}:
\begin{itemize}[leftmargin=1.2em]
\item A \emph{new} stream is introduced by the slot of the edge that delivered its first
transaction; a delivery out of the node's own relay inbox by $\mu$; out-of-band connection
signaling by $\lambda$.
\item A stream leaves the table only when it is spent and holds no relay custody, by
Lemma~\ref{lem:evict}.  Spent streams are retired lazily, by the next admission that reads the
clock past them, and their records are removed in that admission's batch.
\item The record of a stream stores $(\pi(s), x_s+\iota+\sigma)$, so a restart recovers the
charges.
\end{itemize}
Parameters:
\begin{itemize}[leftmargin=1.2em]
\item $f=32$ new streams per slot, $\Pool=3328$, $L=1024$ and $M=S_\mu=768$, so
$B(3)=328\cdot32+3328+1024+768=15\,616$ and $B(s)=B(3)+64(s-3)$;
\item $B(s)\le C=16\,384$, the table's hard capacity, iff $s\le15$: up to 15 successors
Theorem~\ref{theorem:isolation}(1) holds and the floors bind first, and past it the hard capacity
refuses before a floor does;
\item the per-origin quota is unchanged.
\end{itemize}
A new stream past its bound is refused with \texttt{TransactionReplayStreamCapacityExceeded}, as
today.  The refusal now lasts until a stream of the refused slot is spent, at most $H+2\sigma$
(100\,min, about 10\,min for default-lifetime traffic) and $H+H_I+3\sigma$ for application
streams, and it reaches only a slot past its floor and its fair claim.
\end{law}

\begin{proposition}[No permanent lockout]\label{prop:nolockout}
A party with $k$ slots and any number of origins can refuse new streams only through those $k$
slots, and only once the pool is full.  For every other slot, a new origin is refused once that
slot holds $f+\lceil\Pool/n\rceil$ unspent streams.  Every refusal ends within $H+2\sigma$, or
$H+H_I+3\sigma$ for an application stream, of the last admission on the streams that fill the
bound.
\end{proposition}
\begin{proof}
Theorem~\ref{theorem:isolation}, and Lemma~\ref{lem:evict} for the end of a refusal.
\end{proof}

The application horizon is the price of offline mail: a stream must outlive every hold of its
messages.  It lengthens what Proposition~\ref{prop:rotation} costs an honest newcomer in an
application-charged slot to $H_I$, not what any principal can hold.

Inbox mail draws on $\mu$, never on $L$, so a flood of mail cannot refuse this node's
own signaling.  The residual: once mail from $M$ distinct origins is unspent, mail from a new
origin is refused at delivery for up to $H_I+H+3\sigma$; its element stays held, unless its own
hold horizon passes first.  A slot would not do better: a node's inbox mail all enters through
its own store or one member's fetch reply, so charging the delivering slot isolates nothing.

A relay or seed that forwards many distinct origins can fill its own floor legitimately; its
overflow then holds a fair claim of the pool, $\lceil\Pool/n\rceil\ge11$ at worst and $3328$
alone, and past it the honest cost is a refusal on that edge until its streams are spent.  The sender table holds this node's own streams, so it is local.

\section{The removal store, re-keyed}\label{sec:removal}

Law~6 of~\cite{RemovalStability} charged each remove to its sender's origin.  We re-key it
under Law~\ref{law:edge}:
\begin{itemize}[leftmargin=1.2em]
\item an operation or a sync is charged to the slot of the edge that delivered it, and this
node's own writes and repairs to $\lambda$; a read repair of a remote lookup reply is charged to
the slot that delivered the reply, so lookup replies never draw on $L$;
\item a hand-off sync of a placement this node's view makes it own, from the node its view names
as the placement's previous owner (its predecessor or a successor), is charged to the custody
principal $\chi$; every other sync from a predecessor or successor is charged to its slot, as
any edge's.
\end{itemize}
Parameters: $f=128$ removes per slot, $\Pool=16\,384$, $L=7\,168$ and $K=S_\chi=16\,384$, so
$B(3)=328\cdot128+16\,384+7\,168+16\,384=81\,920$ and $B(s)=B(3)+256(s-3)$.  A store files
each item in at most 239 bytes, the widest when every record holds one remove (its key, its
DID, one remove and its charge), so $B(3)$ items take at most about 19.6\,MB.
Theorem~\ref{theorem:isolation} then closes the Sybil residual that \issue{907} was to close: one
mechanism serves both tables.

A store admits one carrier per transition, at most $T_{\max}=1024$ removes, so all-or-none
binds per topic.  The binding limit of a hand-off is cumulative: without $\chi$, every topic of
a range a neighbour hands over draws on that neighbour's one slot, $f+\lceil\Pool/n\rceil$ in
all, and a busy range (about $10$ removes a minute per topic stays $\approx1000$ unstable) would
wait up to $H+\sigma$ while the new owner misses its reads.  The residuals of $\chi$: a
Byzantine predecessor or successor can fill $K$ (a non-goal, as the holder class was), and an
honest range of more than $K$ unstable removes still waits for them to stabilize.

\section{Received frames (\issue{934})}\label{sec:frames}

Let lane $(e,\ell)$ have advertised credit $A$ and released count $R$.  Its liability is
$\liab=A-R$: the frames the receiver may still have to hold, held, in flight, or credited and
not yet used.  Today $\liab\le W=16$ per lane, but every lane holds its whole window, so
$\sum\liab$ is bounded only by the connection count: about 1.3\,GiB at $\Rg(3)$ edges.

\begin{definition}[Reclaim]\label{def:reclaim}
Every lane carries a generation $g$, and its credit frames name it: $\mathsf{Credit}(A,g)$.  Two
frames are added, each on the lane it concerns, never charged against credit, and idempotent
per generation:
\begin{itemize}[leftmargin=1.2em]
\item $\mathsf{Shrink}(L',g)$, receiver to sender: ``from generation $g$, your limit is $L'$'',
with $L'=\min(R+g_\ell,A)$.
\item $\mathsf{Returned}(B,g)$, sender to receiver: ``I committed or reserved $c$ frames, and
from now on I commit at most $B=\max(c,L')$''.
\end{itemize}
The sender takes $B$ as its limit and ignores a credit of an older generation than its own.
Until the
$\mathsf{Returned}$ arrives, the receiver counts the lane's liability at its old $A$ and
advertises nothing on it; after it, $A=B$.
\end{definition}

\begin{law}[Hard frame bound]\label{law:frames}
Every lane has a floor $g_\ell$ and a cap $W$: $g_\ell\le\liab(e,\ell)\le W=16$ always.  The
priority lane's floor is its whole window, $g_{\mathrm{prio}}=W$, and it never borrows; each of
the three data lanes has $g_\ell=1$, so the first frame of a burst on any lane never waits a
round trip.
\begin{itemize}[leftmargin=1.2em]
\item \emph{Demand.}  The admission of a data lane's frame is its demand for pooled credit: a
sender with queued frames always has its floor frame to send, or a frame the receiver holds,
whose admission was the demand.  So no demand frame is needed.
\item \emph{Pool.}  The receiver grants a data lane credit beyond its floor only from the
node's frame pool, at most $W-g_\ell$ to one lane, and to one edge at most its fair claim of
Law~\ref{law:edge}: with $p_e=\sum_\ell(\liab(e,\ell)-g_\ell)^+$ and $n\ge1$ the edges that hold
pooled credit or wait for it,
\begin{gather*}
\textstyle p_e\le\lceil\Pool_F/n\rceil\ \text{when lent},\qquad\sum_e p_e\le\Pool_F,\\
\textstyle\sum\liab\le B_F(s)=\Rg(s)\,(W+3)+\Pool_F .
\end{gather*}
\item Demands are served first come, first served, those of edges below their share
$\lfloor\Pool_F/n\rfloor$ before those of edges between share and claim: a demand at its edge's
claim waits without holding back a later one, and a lane lent more advertises it at once.
\item \emph{Rebalance.}  While a demand of an edge holding less than its share
$\lfloor\Pool_F/n\rfloor$ waits for a full pool, every edge holding more than its share is
shrunk to it.  The shares fit the pool, $n\lfloor\Pool_F/n\rfloor\le\Pool_F$, as the claims need
not; an edge at its share shrinks no other.
\item \emph{Idle.}  A data lane that holds pooled credit and admitted no frame for
$T_{\mathrm{idle}}$ is shrunk to its floor.
\end{itemize}
Parameters:
\begin{itemize}[leftmargin=1.2em]
\item $\Pool_F=256$ frames (16\,MiB, the soft limit it replaces);
\item $T_{\mathrm{idle}}=1$\,s, on a tick the transport is driven by: the wall clock in a
backend, the virtual clock in the deterministic simulation;
\item $B_F(3)=328\cdot19+256=6\,488$ frames, about 405\,MiB at 64\,KiB, and
$B_F(s)=B_F(3)+38(s-3)$ frames.
\end{itemize}
This is independent of how many connections are open.
\end{law}

\textbf{Why $g_{\mathrm{prio}}=W$.}  The most priority-lane frames one edge had outstanding
(held, in flight, or credited and unused) were measured over the dummy suite, with the window
lifted and credit returned per frame: at most $10$ under asynchronous delivery; $16$ in every
sync-storm case; $17$ under a controlled scheduler that holds deliveries on purpose, an
artificial round trip.  Control above $16$ already waits for credit today, and still does; no
smaller floor covers the bursts.  The frame count prices a control frame at the 64\,KiB a data
frame may take: byte-denominated credit (\issue{937}) removes that artifact, and with it most
of $\Rg(s)\,W$.

\textbf{The cost of a fixed $T_{\mathrm{idle}}$.}  The transport has no round-trip estimate,
so $T_{\mathrm{idle}}$ is a constant.  On a link whose round trip exceeds 1\,s, a lane that
pauses for one round trip is shrunk, and its next burst waits one round trip for pooled credit
again; no frame is refused.

\begin{proposition}[Frames]\label{prop:frames}
Under Law~\ref{law:frames}, for an honest sender:
\begin{enumerate}[leftmargin=1.5em]
\item $\sum\liab\le B_F$ and $g_\ell\le\liab(e,\ell)\le W$ in every interleaving, with credit,
shrink and return frames reordered and duplicated;
\item no honest frame is refused: a frame sent under credit is within the liability counted for
it;
\item a lane with queued frames is never blocked forever while the receiver keeps consuming
(progress);
\item an idle lane holds at most its floor after $T_{\mathrm{idle}}$ plus one round trip, so it
starves no active lane, and its first frame after the pause is sent at once;
\item \emph{throughput}: an edge whose data lanes have frames queued holds, within one round
trip of its demand, or two when the pool must be rebalanced first, at least
$\min\bigl(3(W-1),\lfloor\Pool_F/n\rfloor\bigr)$ pooled frames besides its floors, $n\le\Rg(s)$ the
edges holding or waiting for pooled credit, however busy the others are;
\item control traffic never waits for the pool.
\end{enumerate}
\end{proposition}
\begin{proof}[Proof sketch]
(1)~Pooled credit is granted only when the counted liability fits, and the counted liability
never falls below the true one: a shrink takes effect only once its $\mathsf{Returned}$ confirms
what was committed, and a credit of an older generation is ignored.
(2)~A frame committed under $A$ is counted until released, and $B\ge c$.
(3)~A blocked lane's last admission is a pending demand.  The pool refills as frames are
released and idle lanes are shrunk, and service is first come, first served.
(4)~The shrink frees the lane's pooled credit, and its floor stays.
(5)~A waiting edge below its share is lent from what is free; when nothing is, every edge above
its share is shrunk to it, and the frames above the shares return within one round trip, which
leaves the waiting edge its share, since $\sum_{e'\ne e}\min(p_{e'},\lfloor\Pool_F/n\rfloor)\le
(n-1)\lfloor\Pool_F/n\rfloor\le\Pool_F-\lfloor\Pool_F/n\rfloor$.
(6)~The priority lane's floor is its whole window, reserved per edge.
These are the model-checking obligations of the implementation, over the credit model with
shrink and return frames.
\end{proof}

The guarantee~(5) is the price of a hard bound: before \issue{934} every lane had $W$ frames
per round trip whatever the others did.  Now $n$ busy edges share $\Pool_F$, so each keeps its
whole window on all three data lanes while $n\le5$ ($3(W-1)=45\le\lfloor256/5\rfloor=51$), and
the floors plus $\lfloor256/n\rfloor$ pooled frames per round trip beyond; past 256 busy edges,
the floors alone.  Without the
claim, six busy edges would hold the whole pool, and every other data lane would be held to its
floor, one frame per round trip, for as long as they keep sending.

\textbf{Why the model shrinks once.}  The model of Proposition~\ref{prop:frames} lets a lane be
shrunk once per path, which keeps its state space finite.  One is enough: the transitions read
a generation only to compare it with another, so a state of generation $g$ whose older
signals have all been ignored or delivered is, renamed by $g\mapsto0$, a state the model
reaches with no shrink; and a signal of an older generation is a no-op at either end.  So a
second shrink starts from a state the first already covers, and one shrink with credits of the
older generation still in flight, delivered, duplicated and reordered across it, covers them
all.

\textbf{Composition.}  Received frames and the mailbox are disjoint memories: a frame counts in
$B_F$ until it is handed to the mailbox, which happens only once its mailbox charge, and for a
first chunk its item of $O$ (Def.~\ref{def:open}), is granted.  So inbound memory is at most
$B_F\cdot64\,\mathrm{KiB}+M$, and neither bound consumes the other.  A first chunk that waits
for $O$ holds its frame and its credit: backpressure, counted in $B_F$, whose pooled part the
idle shrink returns.  The per-lane cap keeps Proposition~\ref{prop:noref}'s grace $Wd_0$ sound
under pooled credit.

This is a wire change: a total cutover, which wire version 2 names (\issue{1104}).

\section{Reassembled messages (\issue{932})}\label{sec:reassembly}

A message of $\mathit{len}$ bytes travels as chunks $[i,n]$ on one lane $(e,\ell)$, its
\emph{carrier}, which is its class's lane: the sender cuts chunks of equal length $d_0$, the
last shorter, and keeps them contiguous on the carrier.  Let $r(x)=2x$ be the retained bytes of
$x$ received bytes (the buffer and the completed copy), $W$ the lane window
(Law~\ref{law:frames}), $\mathit{TTL}$ the largest time to live a chunk may carry, and
\[
\rho=\frac{\texttt{max\_message\_bytes}}{\mathit{TTL}}
=\frac{60\,\mathrm{MB}}{100\,\mathrm{min}}=10\ \mathrm{B/ms}.
\]

\begin{definition}[Open messages]\label{def:open}
The table $O$ holds, in bytes, the reservation $r(n\,d_0)$ of every message whose first chunk
was admitted and that has neither completed nor expired, charged to the slot of the edge that
delivered that chunk (Law~\ref{law:edge}, an item being a byte), with
$f_O=256$\,KiB, $\Pool_O=128$\,MiB, $L_O=0$ and no dedicated share:
\[
B_O(s)=\Rg(s)\,f_O+\Pool_O,\qquad B_O(3)=210\ \mathrm{MiB},
\]
within the mailbox's $M=256$\,MiB for $\Rg(s)\le512$, that is $s\le94$.  An item of $O$ is
a byte of the mailbox too: $O$ bounds what open messages take of $M$, and $M-B_O$ stays for
every other arrival.
\end{definition}

\begin{definition}[Credit-time]\label{def:credittime}
Let an open message be sampled at instants $t_0<t_1<\dots$: its first chunk's admission, each
later chunk's, and the receiver's sweeps, at most $\Delta$ apart.  At $t_k$ let $c_k=A-N$ be the
credit on the carrier that the sender holds unused: advertised, $A$, and not yet used by an
arrival, $N$.  Its credit-time advances over $[t_k,t_{k+1})$ by
\[
\hat\tau\mathrel{+}=\min\!\bigl(t_{k+1}-t_k,\;c_k\,d_0/\rho\bigr),
\]
the time its unused credit could have taken to send at $\rho$.  Credit the receiver holds back,
by holding the frames that took it, batching its release, or not lending it pooled credit
(Law~\ref{law:frames}), is not in $c_k$, so $\hat\tau$ pauses while the receiver throttles the
sender.
Read at the receiver, $c_k$ also counts credit still in transit to the sender, and frames the
sender already spent it on that are still in flight; neither is the sender's doing, and the rate
floor's grace covers both (Prop.~\ref{prop:noref}).
Let $b$ be the message's bytes received.
\end{definition}

\begin{law}[One charge per message]\label{law:onecharge}
The first chunk $[0,n]$, of length $d_0$, reserves $m=r(n\,d_0)$ once, as an \emph{arrival}
that waits in resource order and is never refused for capacity:
\begin{itemize}[leftmargin=1.2em]
\item on the mailbox ledger, on the edge's budget, on the share of its carrier's class lane,
and on the shared pool; and
\item in $O$, under Law~\ref{law:edge}: within the slot's floor, or else within the pool and
its fair claim, as a whole.
\end{itemize}
It is valid only while $(n-1)\,d_0<\texttt{max\_message\_bytes}$: an honest message of at most
that many bytes has $n=\lceil\mathit{len}/d_0\rceil$, so $n\,d_0$ passes its length by less than
one chunk.
\begin{itemize}[leftmargin=1.2em]
\item Later chunks must have length $\le d_0$, and they charge nothing: they wait only behind
the earlier arrivals of their own stream, which is ordered per peer, lane and carrier, so behind
their own message's chunks alone.
\item The reassembly buffer and the completed message live inside the one reservation.  At
completion the message leaves $O$ and its mailbox reservation moves with it into dispatch, on
its class lane: a move, not a new charge.
\item \emph{Rate floor.}  The message expires, attributably, once
$b<\rho\,\hat\tau-W d_0$.  It also expires at its own $\mathit{ts}+\mathit{ttl}$.  Either
expiry releases the reservation.
\end{itemize}
The separate reassembly byte budget and its refusal ($\mathsf{Rejected}(\mathsf{Capacity})$)
are retired.  The id and slot caps stay, as validity bounds.
\end{law}

\begin{proposition}[No honest message refused]\label{prop:noref}
Under Law~\ref{law:onecharge}, no honest chunked message is refused or dropped for capacity, no
wait can deadlock, and a sender whose carrier delivers at least $\rho$ while it holds credit is
never expired by the rate floor.
\end{proposition}
\begin{proof}
The only charge is the first chunk's, which waits as an arrival before it holds a lane ticket.
The ordering laws of \texttt{ResourceOrderedQueue} cover it, and the capacity it waits for is
released by events that do not wait on it: a later chunk waits for no capacity, and its stream
holds no other message's first chunk, so no message's completion waits on capacity another
message holds.  The transition that used to re-charge, and could only be refused, now moves
bytes already reserved.

Rate floor: over $[t_k,t_{k+1})$ the sender holds $c_k$ unused credits and a backlog (its
message is contiguous on the carrier), so it delivers $\min(c_k d_0,\rho(t_{k+1}-t_k))$ in it,
which is what $\hat\tau$ advanced, times $\rho$, once its credit has reached it.  Between the two,
$\rho\hat\tau$ runs ahead of $b$ only by the lane's outstanding liability: the frames the
receiver still holds, the frames in flight to it, and the credit in transit to the sender, each
counted in $c_k$ or missing from $b$.  That liability is at most $W$ frames, the lane window
(Sec.~\ref{sec:frames}), however long the round trip.  So $b\ge\rho\hat\tau-W d_0$ throughout:
the grace covers every outstanding credit, and a high-latency honest link is never expired.
\end{proof}

\begin{proposition}[Open messages]\label{prop:open}
Let $G=(W+1)\,d_0/\rho+\Delta$, at most $107$\,s plus one sweep for $d_0\le64$\,KiB.  Under
Law~\ref{law:onecharge}:
\begin{enumerate}[leftmargin=1.5em]
\item a first chunk through slot $\kappa$ with $\use_\kappa+m\le f_O$, and every other arrival,
never waits for $O$ behind another edge, however many open messages other edges hold; within
its own slot, first chunks take the floor first come first served
(Def.~\ref{def:resorder}); open messages hold at most $B_O$ of the mailbox;
\item an honest message that waits for the pool of $O$ waits at most $G$ beyond the honest
messages ahead of it: an open message that sends nothing is expired within $G$, and a slot's
next first chunk queues in resource order behind the claim already waiting;
\item past $G$, every open message has $\rho\,\hat\tau\le b+W d_0$: the credit-time for which it
holds its reservation is paid for by bytes its sender streamed.
\end{enumerate}
\end{proposition}
\begin{proof}
(1)~Theorem~\ref{theorem:isolation}(2) for $O$; an arrival that is not a first chunk takes no
item of $O$.  (2)~A sender that holds its credit and sends nothing has $\hat\tau$ grow at the
rate of time and $b=d_0$, so it meets $d_0<\rho\hat\tau-Wd_0$ after $(W+1)d_0/\rho$, and the next
sweep expires it; resource order serves the waiting claim before a later one.
(3)~The rate floor.
\end{proof}

\textbf{Residual.}  The receiver cannot tell a sender that withholds a message from one whose
scheduler is busy with higher-priority classes between chunks: both hold credit and send
nothing.  So a large message on a link that leaves its class less than $\rho$ is expired.  This
is the cost of the rate floor; $\rho=\texttt{max\_message\_bytes}/\mathit{TTL}$ is the least
rate at which holding a reservation is no cheaper than before \issue{932}, when it cost the
bytes held.  A message so expired is the sender's to send again.

\section{Hot paths (\issue{938})}\label{sec:hot}

These are performance obligations, kept apart from the security laws above.

\begin{definition}[Resource order]\label{def:resorder}
A request $w$ has a stream $s(w)$, a footprint $F(w)$, the resources it would draw on now, and
reliefs $\Rg(w)=\{(\rho,D)\}$: a resource $\rho\in F(w)$ whose release may shed $D\subseteq
F(w)$.  The mailbox has $(\mathit{Lane}(\ell),\{\mathit{Shared}\})$ while the reservation of
$\ell$ does not cover $w$, and $(\mathit{Floor}(\kappa),\{\mathit{OpenPool}\})$ while a first
chunk is beyond the floor of its slot $\kappa$ but within it once $\kappa$ holds nothing.  A
waiter refused by the ledger records its refusal $R(w)$, one budget's resources.  With
$B(w)$ the blockers of $w$ below,
\[
\begin{aligned}
B(w)&=\{s(w)\mid\exists v<w.\;s(v)=s(w)\}\cup\{r\in F(w)\mid\exists v<w.\;r\in H(v)\},\\
H(v)&=\begin{cases}
R(v) & v\ \text{refused},\\
\{\rho\mid(\rho,D)\in\Rg(v)\} & B(v)\ne\emptyset,\
  B(v)\subseteq\textstyle\bigcup_{(\rho,D)\in\Rg(v)}\{\rho\}\cup D,\\
\emptyset & \text{otherwise},
\end{cases}
\end{aligned}
\]
and $w$ waits iff $B(w)\ne\emptyset$, else it is attempted.  A waiter blocked only for lack of
room that a relief may make \emph{holds} the relief; one blocked by its stream or by its own
budget holds nothing.
\end{definition}

\begin{proposition}[Order on reservations]\label{prop:resorder}
Under Def.~\ref{def:resorder}:
\begin{enumerate}[leftmargin=1.5em]
\item a frame within its lane's reservation never waits behind a borrower of $\mathit{Shared}$,
nor behind another peer's budget;
\item a lane's reservation, and a slot's floor in $O$, are served first come first served: a
later, smaller request does not pass an earlier one waiting for the room on them.
\end{enumerate}
\end{proposition}
\begin{proof}
(1)~Within its reservation, $\mathit{Shared}\notin F(w)$; a refusal by the node budget is
$\{\mathit{Lane}(\ell'),\mathit{Shared}\}$, met only on $\ell'=\ell$, which is a borrower of
$\ell$'s reservation room, not of $\mathit{Shared}$; another peer's budget
$\mathit{Peer}(p')\notin F(w)$, and a waiter blocked by it holds nothing.  (2)~A waiter
holding $\mathit{Lane}(\ell)$ blocks every later $w$ with $\mathit{Lane}(\ell)\in F(w)$, so
every later request of $\ell$; likewise $\mathit{Floor}(\kappa)$.
\end{proof}

Item~2 is a change: before \issue{938}, a blocked waiter held nothing, so smaller frames could
take a reservation's room ahead of a larger one indefinitely.  It is the first law of resource
order, no overtaking on an exhausted resource, applied to the room a reservation frees.

\begin{obligation}[Admission complexity]\label{ob:complexity}
\begin{itemize}[leftmargin=1.2em]
\item An arrival, release or cancel in \texttt{ResourceOrderedQueue} costs
$O((g+k+r)\log n)$ in the $n$ waiters, where $g$ is the requests it grants and $r$ the
resources it names.  A blocked waiter is parked under one blocker and its holds, a refused one
under its refusal.  An event scans each freed resource's parked waiters in arrival order.  It
stops at the first that parks under it again, since every later one is blocked behind that one;
an invariant, one refusal per resource, keeps this rule exact.  The term $k$ counts the waiters
a scan judges and parks under a different blocker: a resource the event freed was the blocker
they were indexed by.  By the stop rule, a scan judges at most one waiter beyond its grants and
its $k$.  Each such move is charged to the freeing of the blocker it leaves, at most once per
freeing.
\item A received frame takes its link's lock once to be admitted and once to be released.  It
takes the node pool's lock only while its data lane wants pooled credit or its charge changes,
at most once per lend, shrink or return of the lane (Law~\ref{law:frames}), never per frame
in steady state.
\item The pool's service costs $O(\log n)$ in the $n$ edges per demand it serves: the
claimants are counted incrementally, and demands are indexed, so a service visits only the
demands it can serve.
\item Each peer runs one admission watcher, not one task per send.
\item Unread admissions are indexed by peer.
\end{itemize}
They are verified on the sync-storm matrix by operation counts, not by wall time, which the
deterministic simulation does not model: lock acquisitions per received frame, waiters visited
per queue event, demands visited per service, tasks spawned per send, and inbox items scanned
per withdrawal, counted by test-only counters on the same seeds before and after, and each
bounded by a regression assertion at the matrix's largest size.
\end{obligation}

\section{The reassembly barrier (\issue{927})}\label{sec:barrier}

The inbound actor of a connection has one lane per category $\kappa\in K=\{\mathit{ctl},
\mathit{sto},\mathit{e2e},\mathit{app}\}$ and a reassembly lane $\rho$, each serving one event
at a time.  Events take sequence numbers $\mathit{seq}$ in ticket order.  A reassembled message
keeps its final chunk's sequence and enters the lane of its category.  Before \issue{927} a data
lane waited behind every earlier unfinished reassembly, whatever its category, and the control
lane waited behind none: the \emph{control exemption}.

Two facts of the code refute the old argument (``at most one frame in hand-over per lane, and
control is never chunked'').
\begin{enumerate}[leftmargin=1.5em]
\item \emph{Pacing.}  An arrived frame is submitted awaited: its carrier's drainer hands over the
next frame only once this one completed, and a final chunk's completion moves onto its
reassembled message.  So on the paced path each carrier has at most one event in the actor.
Frames released from the pre-admission hold are not paced: up to the peer's event share,
$32$, enter at once, at admission.  Behind such a burst a control frame without the exemption
would wait for up to $32$ reassembly frames.
\item \emph{Chunked control.}  A payload is chunked whenever it exceeds the peer's advertised
\texttt{max-message-size}, which has no floor (\issue{1130}), so control messages carrying an
SDP can arrive chunked, on the control carrier.  Under the exemption, a control frame released
after them in the same burst ran before their reassembled message: control reordered.
\end{enumerate}

\begin{definition}[Per-category barrier]\label{def:barrier}
The category $\kappa(e)$ of an event is its lane's for a whole frame, and its carrier's for a
chunk; a chunk with no carrier, which only a test hands in, has $\kappa(e)=\top$.  Let $U$ be
the reassembly events ticketed and not yet completed.  A lane $\kappa$ dispatches its front $f$
iff
\[
\neg\exists u\in U.\;\mathit{seq}(u)<\mathit{seq}(f)\wedge\kappa(u)\in\{\kappa,\top\}.
\]
\end{definition}

\begin{proposition}[Order per category]\label{prop:catorder}
Under Def.~\ref{def:barrier}, the messages of one category are processed in sequence order,
whole or chunked.  A frame of category $\kappa$ waits only behind unfinished reassembly of
$\kappa$: at most one frame on the paced path, and at most the $\kappa$-chunks of one released
burst.
\end{proposition}
\begin{proof}
A lane serves in sequence order.  A reassembled message $m$ of $\kappa$ has the sequence of its
final chunk $u$; every $\kappa$-frame $f$ with $\mathit{seq}(f)>\mathit{seq}(u)$ waits while
$u\in U$, and $u$ leaves $U$ in the step that inserts $m$ into lane $\kappa$ at
$\mathit{seq}(u)$, ahead of $f$.  The lane-mismatch refusal makes $m$'s category its carrier's,
so $\kappa(u)$ is $m$'s lane.  The bounds are pacing's and the hold's.
\end{proof}

\textbf{Order is guaranteed per category only, and no protocol relies on cross-category
order.}  Categories travel on separate data channels (\issue{906}), whose arrival order the
transport does not relate, and the actor serves its lanes concurrently, so whole frames of two
categories were never ordered; the old barrier ordered data only behind earlier reassembly.
Every request and its report, and every exchange whose steps must be ordered, lies in one
category:
\begin{center}
\footnotesize
\begin{tabular}{@{}p{0.36\columnwidth}p{0.56\columnwidth}@{}}
\toprule
Pair & Why order across categories is not needed\\
\midrule
storage sync and its report (\textit{sto}) against topology control (\textit{ctl}) & the
sync names its signed destination and is forwarded by the current topology when not ours;
persistence and acks are keyed by entry\\
relay custody, holds and their acks against inbox fetches & all \textit{sto}; relayed mail
and a direct \texttt{CustomMessage} take different paths and were never ordered\\
link control against payloads & handled before the actor; keyed by delegation digest, with
lease durations; an \texttt{Unknown} means the referencing frame was dropped and is resent\\
E2E handshake and stream against custom messages & the handshake and stream are all
\textit{e2e}; a custom message needs no E2E session\\
DHT control requests and reports; storage requests and reports & each pair within
\textit{ctl}, or within \textit{sto}\\
\bottomrule
\end{tabular}
\end{center}

\textbf{Decision.}  Replace the node-wide barrier and its hand-off barrier by
Def.~\ref{def:barrier}.  The control exemption is gone as a special case: control waits only
behind chunked control, which a peer causes only by advertising a limit below its control
messages.  A data category no longer waits behind another category's reassembly.  The ablation
toggle \texttt{ProtectionLayer::BarrierControlExemption} and its legacy feedback-loop witness
are removed; the storm witnesses that a control probe behind another category's reassembly
backlog is not blocked, and a regression test that a chunked control message released from
the hold keeps its place among later control.

\section{Summary of bounds}\label{sec:summary}

\begin{center}
\footnotesize
\begin{tabular}{@{}p{0.22\columnwidth}p{0.3\columnwidth}p{0.36\columnwidth}@{}}
\toprule
Table & Bound & At the bound\\
\midrule
frames held & node $B_F(s)$; lane floor and cap $W$ & sender waits for credit (backpressure)\\
mailbox bytes & edge budget, lane share, pool & arrival waits (backpressure)\\
open messages & $\Rg(s) f_O+\Pool_O$ by slot, fair pool, in the mailbox & first chunk waits
(backpressure); expired below $\rho$\\
replay streams & $\Rg(s) f+\Pool+L+M$ by slot, fair pool & new stream refused, $\le H+2\sigma$
($H+H_I+3\sigma$ application)\\
removal store & $\Rg(s) f+\Pool+L+K$ by slot, fair pool & write refused, unacked\\
origin & rate quota within its edge & refused for the rate window\\
application inbox & 64\,MiB & dropped, counted (by design)\\
\bottomrule
\end{tabular}
\end{center}

\textbf{What a free identity consumes.}  An origin alone consumes nothing beyond its quota on the
edge that carries it.  An edge consumes its slot's floors, which outlive it until its items
leave.  The pool is shared, and is the most that any number of free identities can take
together.

\section{Implementation order}\label{sec:plan}

\begin{center}
\footnotesize
\begin{tabular}{@{}llp{0.6\columnwidth}@{}}
\toprule
PR & Issue & Law\\
\midrule
1 & \issue{907} & Laws~\ref{law:edge}, \ref{law:streams}; Sec.~\ref{sec:removal}\\
2 & \issue{932} & Law~\ref{law:onecharge}\\
3 & \issue{934} & Law~\ref{law:frames} (wire cutover)\\
4 & \issue{938} & Obligation~\ref{ob:complexity}, measured\\
5 & \issue{927} & Sec.~\ref{sec:barrier} (subtraction)\\
6 & \issue{1094} & Def.~\ref{def:principals}; Prop.~\ref{prop:rotation}\\
\bottomrule
\end{tabular}
\end{center}

PR~1 adds the slot-accounting primitive that PR~2 reuses for the mailbox, and the relay
element horizon of Definition~\ref{def:relayhorizon}.  PR~4 rewrites the
queue that PR~2 changes, so it follows PR~2.  PR~5 and PR~6 need nothing from the others.

\textbf{Operator steps}:
\begin{itemize}[leftmargin=1.2em]
\item PR~1 changes the replay record, which gains its introducer and spent instant.  Remove the
\texttt{transaction-replay} store while the node is stopped; this is the documented reset of
the durable-records law.  The removal store gains a slot charge, so remove
\texttt{dht-removals} too.  Held mail now expires $H_I$ after its hold, even in a carrier that
newer holds keep alive.
\item PR~3 is a wire cutover with no step.
\end{itemize}

\sloppy
\bibliographystyle{IEEEtran}
\bibliography{load-bounds}

\end{document}
